% Paper B (problem131_multistate.tex): material removed from the paper on
% 2026-09-26 while trimming the appendix. Every piece below is verbatim from
% the file and lines named in its header, before the trim. Labels and
% cross-references refer to that earlier version, so this file is not meant
% to be compiled.
%
% What replaced each piece in the paper:
%  1-2.  Lemma start-weights(3) became Remark rem:start-weights (proof omitted).
%        Its proof and the 2 integral forms it used are below.
%  3.    Corollary cor:start-modes became Remark rem:start-modes, which cites
%        Paper C's face-maxima and tie-window theorems. Its old proof is below.
%  4.    Example ex:K3-start became a short remark of the same label, and
%        Remark rem:K3-finite became a short remark in general_start.tex. The
%        long remark and Table tab:K3-window are below.
%  5.    Theorem thm:frontier-superlevel became Remark rem:frontier-superlevel
%        in frontier_modes.tex (proof omitted). Its proof is below.
%  6.    The covariance remark rem:tetra was folded into the proof of
%        Proposition prop:field-tilt(3) (variance only).
%  7-10. Lemmas lem:simplex-integral and lem:simplex-laplace, the first
%        paragraph of the proof of the asymptotic assertions, and the first
%        part of the proof of Theorem thm:simplex-interior were replaced by
%        Lemma lem:no-rare, which reads the proof of the crossover theorem with
%        no rare counts. The moment bound in the proof of Lemma
%        lem:start-weights(2), which used lem:simplex-integral, was reproved
%        with the integration by parts of the majorant proof. The old versions
%        are below.
%  11.   Remark rem:last-exit-trees became one sentence in the related work.
%  12.   Lemma lem:singleton-insertion (not used by any proof) was removed.
%  13-14. The near-vertex side results (Theorem thm:simplex-fixed-vertex, the
%        all-singleton case, Corollary cor:vertex-shape-law and Proposition
%        prop:vertex-probability-failure) were removed, and one sentence in
%        frontier_boundary.tex mentions them (proofs omitted).
%  15.   Sentences in the main file that pointed to removed results.

%======================================================================
% 1. Lemma start-weights(3): statement and the paragraph after the lemma
% Source: general_start.tex, lines 27-39 of the version before 2026-09-26 trim.
%======================================================================
\item Fix a compact set $\mathcal K$ in the open face and $0<c<C$, and put
$\alpha=n/N$ and $A=\sum_{b\in F}\sqrt{\alpha_b}$. Uniformly for
$\alpha\in\mathcal K$ and $cL\le T\le CL$,
$\omega_a(n)=\sqrt{\alpha_a}/A+O(1/L)$.
\end{enumerate}
\end{lemma}

Under the uniform start, $\omega_a(n)$ is the probability that the chain starts
in $a$ given $K_N=n$. So by (3) a path conditioned to have fractions $\alpha$
starts in $a$ with probability close to $\sqrt{\alpha_a}/A$, and not $\alpha_a$.
In the notation of Proposition~\ref{prop:kirchhoff}, $\sqrt\alpha$ is the Perron
vector $h$ of the tilted generator. The proof is in
Appendix~\ref{app:start-field}.

%======================================================================
% 2. Appendix B: the integral forms eq:startw-integral and eq:startw-nu, used only in the proof of (3)
% Source: app_start_field.tex, lines 6-31 of the version before 2026-09-26 trim.
%======================================================================
We prove Lemma~\ref{lem:start-weights}. Throughout $1/d<p<1$, so $\sigma=p-r>0$
and $v=r/\sigma=u/(1-u)$, and $B_m$ is the polynomial of
Lemma~\ref{lem:uniform-laguerre-bessel}. First we write the law in 2 more forms.
The refresh form in Proposition~\ref{prop:dpmf}(a) is linear in $\mu$, and with
$\mu=\delta_a$ it gives $p^{N-1}S^{(a)}_n(u)$. So, with sums over
$1\le m_b\le n_b$ for $b\in F$ and $M=\sum_bm_b$,
\begin{equation}\label{eq:startw-weights}
 \omega_a(n)=\E_w\Bigl[\frac{m_a}M\Bigr],\qquad
 w(m)\propto\frac{M!}{\prod_bm_b!}\,v^M\prod_{b\in F}\binom{n_b-1}{m_b-1}.
\end{equation}
Next write $(M-1)!=\int_0^\infty e^{-t}t^{M-1}dt$ in the refresh form and use
$\sum_m m\binom{n-1}{m-1}y^m/m!=yB_n'(y)$. All terms are positive, so Tonelli's
theorem gives
\begin{equation}\label{eq:startw-integral}
 P^\mu_N(n)=\sigma^{N-1}\int_0^\infty e^{-t}\sum_{a\in F}\mu_a
 B_{n_a}'(vt)\prod_{b\ne a}B_{n_b}(vt)\,dt.
\end{equation}
For the uniform start an integration by parts turns this
into~\eqref{eq:simplex-balance-integral}. Put $\kappa_m(y)=yB_m'(y)/B_m(y)$. Then
$B_{n_a}'\prod_{b\ne a}B_{n_b}=(\kappa_{n_a}/y)\prod_bB_{n_b}$, and
\eqref{eq:startw-integral} gives
\begin{equation}\label{eq:startw-nu}
 \omega_a(n)=\E_\nu\Bigl[\frac{\kappa_{n_a}(vt)}{\sum_b\kappa_{n_b}(vt)}\Bigr],
 \qquad
 \nu(dt)\propto e^{-t}\prod_bB_{n_b}(vt)\,\frac{\sum_b\kappa_{n_b}(vt)}{vt}\,dt.
\end{equation}

%======================================================================
% 2 (cont). Proof of Lemma start-weights(3)
% Source: app_start_field.tex, lines 74-111 of the version before 2026-09-26 trim.
%======================================================================
(3) Let $\alpha_b\ge\varepsilon>0$ on $\mathcal K$, put $Y=Nv$, and note that
$Y=T(1+O(L/N))$. The proof of Lemma~\ref{lem:vertex-elasticity} gives
$\kappa_m(y)=\E[J+1]\ge1$, with $J$ as in that proof, $\kappa_m\le1+\sqrt{my}$ and
$\kappa_m^2+y\kappa_m\ge my+1$ for all $y>0$. The last inequality gives
$\kappa_m\ge\sqrt{my+y^2/4}-y/2\ge\sqrt{my}-y/2$, so
$|\kappa_m(y)-\sqrt{my}|\le1+y$. Since $\sqrt{n_by}=\sqrt{\alpha_b}\sqrt{Ny}$,
the ratio $\rho_a=\kappa_{n_a}/\sum_b\kappa_{n_b}$ satisfies
\[
 \Bigl|\rho_a-\frac{\sqrt{\alpha_a}}A\Bigr|
 \le\frac{2k(1+y)}{A\sqrt{Ny}-k(1+y)}
 \qquad\text{when }A\sqrt{Ny}>k(1+y).
\]
Let $I=[A^2Y/4,4A^2Y]$. For $t\in I$ and $y=vt$ we have $Ny=Yt\ge A^2Y^2/4$ and
$y\le4kY^2/N$, so $|\rho_a-\sqrt{\alpha_a}/A|=O(1/L)$ on $I$. Off $I$ we use
$|\rho_a-\sqrt{\alpha_a}/A|\le1$ and show that $\nu$ gives $I^c$ mass $o(1/L)$.

For the upper bound, Lemma~\ref{lem:uniform-laguerre-bessel} and
$\Phi(x)\le e^{2\sqrt x}$ give $B_m(y)\le y\Phi(my)\le ye^{2\sqrt{my}}$, and
$\sum_b\kappa_{n_b}\le k+A\sqrt{Yt}$. With $\sum_b\sqrt{n_bvt}=A\sqrt{Yt}$, the
density of $\nu$ before normalizing is at most
\[
 (vt)^{k-1}(k+A\sqrt{Yt})\exp\bigl(A^2Y-(\sqrt t-A\sqrt Y)^2\bigr).
\]
Off $I$ we have $(\sqrt t-A\sqrt Y)^2\ge A^2Y/4$. Splitting the square in half
and integrating the other half shows that the mass of $I^c$ before normalizing
is at most $C_3v^{k-1}Y^{C_4}e^{A^2Y-A^2Y/8}$. For the lower bound take
$|\sqrt t-A\sqrt Y|\le1$, an interval of length about $4A\sqrt Y$ inside $I$.
There $y=O(L^2/N)$ and $\alpha_bYt\ge1$ for large $N$. By
Lemma~\ref{lem:uniform-laguerre-bessel}, $B_m(y)\ge y\Phi(my)(1-y/2)$, and the
expansion of $I_1$ in Section~\ref{sec:multistate-model} gives
$\Phi(x)\ge c_0x^{-3/4}e^{2\sqrt x}$ for $x\ge1$, since $\Phi$ is positive and
continuous. With $\sum_b\kappa_{n_b}\ge k$ the density is at least
$c_1v^{k-1}Y^{-C_5}e^{A^2Y-1}$ on this interval, so the normalizing constant is
at least $c_2v^{k-1}Y^{-C_5}e^{A^2Y}$. As $A\ge1$, the mass of $I^c$ is
$O(Y^{C_4+C_5}e^{-Y/8})=o(1/L)$. So $\omega_a(n)=\E_\nu\rho_a
=\sqrt{\alpha_a}/A+O(1/L)$, and all constants depend only on $k$,
$\varepsilon$, $c$ and $C$.
\end{proof}

%======================================================================
% 3. Corollary cor:start-modes (statement and the paragraph after it)
% Source: general_start.tex, lines 118-144 of the version before 2026-09-26 trim.
%======================================================================
\[
 \ell^\mu_1(t)=\log\mu_{(1)},\qquad
 \ell^\mu_k(t)=\log\frac{M_k}k+(k-1)(t-a_k)\quad(2\le k\le d).
\]
\begin{enumerate}[(1)]
\item For every support $S$ with $|S|=k\ge2$ and $\mu(S)>0$,
$\max\{P^\mu_N(n)/p^{N-1}:\supp n=S\}\to\frac{\mu(S)}ke^{(k-1)(t-a_k)}$,
uniformly on compact sets of $t$.
\item Let $J$ be a compact set of $t$ on which $\max_k\ell^\mu_k(t)$ is attained
only at $k=k^*$. For all large $N$ and all $t\in J$, every global mode of
$P^\mu_N$ has exactly $k^*$ positive coordinates, and its support $S$ has
$\mu(S)=M_{k^*}$. The winning size is nondecreasing in $t$.
\item Let $t^*=a_d+\frac1{d-1}\log(d\mu_{(1)})$, where $\ell^\mu_1$ and
$\ell^\mu_d$ cross. No size $2\le k\le d-1$ is the unique winner on an open
interval of $t$ if and only if
$\log\frac{M_k}{k\mu_{(1)}}+(k-1)(t^*-a_k)\le0$ for every $2\le k\le d-1$.
\end{enumerate}
For the uniform start the left side in (3) is $(k-1)(a_d-a_k)<0$, and (2) gives
Corollary~\ref{cor:simplex-modes} in the window.
\end{corollary}

Part (1) follows from Lemma~\ref{lem:confinement} and
Lemma~\ref{lem:start-weights}(3), and (2) and (3) are facts about the upper
envelope of the lines. The proof is in Appendix~\ref{app:start-field}.

\begin{example}[the start decides which face wins]\label{ex:K3-start}
Let $d=3$ and $\mu=(\frac12,\frac12,0)$. The lines are $\log\frac12$ for the

%======================================================================
% 3 (cont). Proof of Corollary cor:start-modes
% Source: app_start_field.tex, lines 216-236 of the version before 2026-09-26 trim.
%======================================================================
\begin{proof}[Proof of Corollary~\ref{cor:start-modes}]
(1) The balanced bins give the lower bound, by
Theorems~\ref{thm:start-crossings}(2) and~\ref{thm:simplex-window}. For the upper
bound put $y=\sqrt L\,(n/N-\frac1k\mathbf 1)$ and fix $R$ with
$e^{-k^2R^2/4}\le1/(2k)$. If $|y|>R$, then $\sum_a\mu_a\omega_a(n)\le\mu(S)$,
and Lemma~\ref{lem:confinement}(b) bounds $P^\mu_N(n)/p^{N-1}$ by
$\frac{\mu(S)}{2k}e^{(k-1)(t-a_k)+o(1)}$. If $|y|\le R$,
Lemma~\ref{lem:start-weights}(3) gives
$\sum_a\mu_a\omega_a(n)=\mu(S)/k+O(R/\sqrt L)$, and $S_n(u)$ is at most its
value at a balanced bin by Theorem~\ref{thm:simplex-balance}.
(2) A vertex gives exactly $\mu_i$, a support with $\mu(S)=0$ gives 0, and among
supports of size $k$ the largest $\mu(S)$ is $M_k$. So by (1) the largest bin
with $k$ positive coordinates, divided by $p^{N-1}$, has logarithm
$\ell^\mu_k(t)+o(1)$, uniformly on $J$, and on $J$ the line $\ell^\mu_{k^*}$ beats
the others by a positive gap. If $k_1$ wins at $t_1$ and $k_2$ at $t_2>t_1$,
adding the 2 comparisons gives $(k_2-k_1)(t_2-t_1)\ge0$.
(3) The function $E=\max(\ell^\mu_1,\ell^\mu_d)$ is convex with its only kink at
$t^*$. A line with slope strictly between 0 and $d-1$ that is at most $E$ at
$t^*$ is below $E$ everywhere else. Otherwise it is above $E$ near $t^*$, and
some proper line is then the unique maximum on an open interval.
\end{proof}

%======================================================================
% 4. Example ex:K3-start (full version) and the paragraph after it
% Source: general_start.tex, lines 146-172 of the version before 2026-09-26 trim.
%======================================================================
$\log\frac14+t-a_2$ for the edges through 3 and $\log\frac13+2(t-a_3)$ for the
full face. Put $t_{\rm ef}=\log\frac32+2a_3-a_2=\frac12\log\frac{32\pi}{243}$. By
Corollary~\ref{cor:start-modes}, and Theorem~\ref{thm:simplex-balance} on the
edge, for every fixed $t\in(a_2,t_{\rm ef})$ and all large $N$ the global modes
are the balanced bins of $\{1,2\}$. This window is $(-0.4673558,-0.4412978)$, of
width $\log\frac{16}{9\sqrt3}=0.0260580$. Below it the modes are $(N,0,0)$ and
$(0,N,0)$, and above it they have 3 positive counts. With the uniform start no
proper face wins in the window (Corollary~\ref{cor:simplex-modes}). More
generally, by Corollary~\ref{cor:start-modes}(3) an edge wins on an open window
on $K_3$ if and only if $(\mu_{(1)}+\mu_{(2)})^2>\frac{3^{5/2}}8\mu_{(1)}$. Here
$1>0.974$, a thin margin, while the uniform start has $\frac49<0.650$.

The lower end is Paper~A's crossing, exactly and for every $N\ge2$, in the
sense that the balanced bins of the edge overtake the vertices exactly there.
Indeed $\mu$ is constant on $\{1,2\}$, so
$P^\mu_N(n_1,n_2,0)/P^\mu_N(N,0,0)=S_{(n_1,n_2)}(u)$, which is the polynomial of
Paper~A~\cite{paperA} at $u=r/p$. So the crossing is at Paper~A's root $u_N$,
that is at $T=Nu_N/(1+2u_N)$, which differs from Paper~A's $z_N=Nu_N/(1+u_N)$ by
$O(L^2/N)$.
\end{example}

Remark~\ref{rem:K3-finite} in Appendix~\ref{app:start-field} describes the
window at finite $N$. Its second-order ends are predictions that we do not
prove, and Table~\ref{tab:K3-window} compares them with exact computations.


%======================================================================
% 4 (cont). Remark rem:K3-finite and Table tab:K3-window
% Source: app_start_field.tex, lines 113-167 of the version before 2026-09-26 trim.
%======================================================================
\subsection{\texorpdfstring{The window on $K_3$ at finite $N$}{The window on K3 at finite N}}

\begin{remark}[the window at finite $N$, not proved]\label{rem:K3-finite}
Here $z=Nu=T/p$, since the finite-size terms depend on the variable. The gain
of the largest bin on a face over the balanced bins, predicted in Section~\ref{sec:general-start}, is
$\frac12\gamma^{\mathsf T}\Sigma\gamma/z$ in Paper~C's terms, where $\gamma$ is
the gradient of $\log\sum_a\mu_a\sqrt{\alpha_a}$ at the center, along the plane
$\sum_a\alpha_a=1$, and $\Sigma$ is the occupation covariance there. For the full face, in all 3 coordinates,
$\gamma=(\frac14,\frac14,-\frac12)$ and $\Sigma=\frac29(I-\frac13J)$, so the gain
is $\frac1{24z}$. The quadratic model behind it puts the largest bin near
$n_3/N=\frac13-\frac1{9z}$, with $n_1=n_2$ up to parity. With the terms
$-\frac1{8z}$ of the edge and $-\frac2{8z}$ of the full face from
Theorem~\ref{thm:start-crossings}, the term at the upper end becomes
$\frac2{8z}-\frac1{8z}-\frac1{24z}=\frac1{12z}$. For a balanced $k$-face bin a
heuristic calculation, which exact computations support, adds a discrete term
$-\frac{(k-1)^2}2\frac{z^2}N$ to the logarithm
of~\eqref{eq:simplex-uniform-ratio}. So the 2 ends should satisfy
\[
 z-L+\tfrac12\log z=a_2+\frac1{8z}+\frac{z^2}{2N}+\cdots,\qquad
 z-L+\tfrac12\log z=t_{\rm ef}+\frac1{12z}+\frac{3z^2}{2N}+\cdots,
\]
where only $a_2+\frac1{8z}$ is proved (by Theorem~\ref{thm:simplex-crossing}).
Without the $z^2/N$ terms the width in $t$ is $0.0260580-0.0546957/L+\cdots$,
and below $N$ of about $10^4$ the $z^2/N$ terms are larger than the $1/L$ terms.
Table~\ref{tab:K3-window} compares these predictions with exact computations.
\end{remark}

\begin{table}[t]
\centering\small
\begin{tabular}{@{}rcccc@{}}
\toprule
$N$ & width & predicted width & $z(\frac13-\frac{n_3}N)$ &
$z\log(M_{\rm full}/P_{\rm bal})$\\
\midrule
300 & 0.0843 & 0.0802 & 0.139 & --\\
3000 & 0.03236 & 0.03211 & 0.1242 & 0.0453\\
30000 & 0.02231 & 0.02258 & 0.1208 & 0.0442\\
300000 & 0.02141 & 0.02166 & 0.1187 & 0.0436\\
limit & 0.02606 & & $1/9$ & $1/24$\\
\bottomrule
\end{tabular}
\caption{The edge window of Example~\ref{ex:K3-start} in $z=Nu=T/p$, with
$t=z-L+\frac12\log L$ (at $N=300$ the difference from $T$ matters), and the
prediction of Remark~\ref{rem:K3-finite}. The last 2 columns are at the upper
end, for the largest bin $n$ with 3 positive counts, of probability
$M_{\rm full}$, and $P_{\rm bal}=P^\mu_N(\frac N3,\frac N3,\frac N3)$. The limit
of the width is proved, and $1/9$ and $1/24$ are predictions. The ends come from
the refresh formula, confirmed by the run formula resolved by the first state,
and the lower end equals Paper~A's root. The largest bin was found by a line
search along $n_1=n_2$ checked on a neighborhood, so for $N\ge3000$ we assume it
lies on or next to that line. A dynamic program over all bins confirmed the
modes at $N=90$, $300$ and $1002$. The predictions were written before the
computations.}
\label{tab:K3-window}
\end{table}

%======================================================================
% 5. Theorem thm:frontier-superlevel, its proof and the paragraphs around it
% Source: app_start_field.tex, lines 270-328 of the version before 2026-09-26 trim.
%======================================================================
The same profile counts the bins that beat a vertex. Put $h_*=\max_ih_i$,
$\widetilde V_{N,h}=\max_i\widetilde P_{N,h}(Ne_i)$ and, for a support $S$ with
$|S|=k\ge2$,
\[
 B_S(t,h)=\ell_S(t,h)-h_*=(k-1)(t-a_k)+\bar h_S-h_*.
\]
Let $|\mathbb B_{k-1}|=\pi^{(k-1)/2}/\Gamma(\frac{k+1}2)$ be the volume of the
unit ball of $H_S$.

\begin{theorem}[Gaussian superlevel geometry]\label{thm:frontier-superlevel}
Put $T=L-\frac12\log L+t$ and, for a count vector with support exactly $S$,
$y_i=\sqrt L\,(n_i/N-1/k)$ for $i\in S$. Uniformly when $t$, $h$ and $y\in H_S$
range over fixed compact sets,
\begin{equation}\label{eq:frontier-profile}
 \frac{\widetilde P_{N,h}(n)}{\widetilde V_{N,h}}
 =\exp\left(B_S(t,h)-\frac{k^2}{4}\sum_{i\in S}y_i^2+o(1)\right).
\end{equation}
Let $M_{N,S}(t,h)$ be the number of bins with support exactly $S$ whose tilted
probability is strictly greater than $\widetilde V_{N,h}$. Then, uniformly for
$(t,h)$ in compact sets,
\begin{equation}\label{eq:frontier-volume}
 \frac{L^{(k-1)/2}}{N^{k-1}}M_{N,S}(t,h)
 \longrightarrow
 \frac{|\mathbb B_{k-1}|}{\sqrt k}
 \left(\frac{4\,[B_S(t,h)]_+}{k^2}\right)^{(k-1)/2},
 \qquad [x]_+=\max(x,0).
\end{equation}
If $B_S(t,h)<0$, there are no such bins for all large $N$. At $B_S(t,h)=0$ the
normalized count tends to 0, but we do not claim that the set is eventually
empty.
\end{theorem}

\begin{proof}
The field factor relative to the best vertex is $e^{h\cdot\alpha-h_*}$, so
$\log(\widetilde P_{N,h}(n)/\widetilde V_{N,h})=\log S_n(u)+h\cdot\alpha-h_*$. For
$y$ in a compact set $h\cdot\alpha=\bar h_S+o(1)$, and
Lemma~\ref{lem:confinement}(a) gives~\eqref{eq:frontier-profile}. Since
$h\cdot\alpha\le h_*$, Lemma~\ref{lem:confinement}(b) with $R$ so large that
$(k-1)(t-a_k)-k^2R^2/4<-1$ on the compact set shows that every bin above
$\widetilde V_{N,h}$ has $|y|\le R$ for large $N$. On this ball the profile
converges uniformly. So for every $\delta>0$ and large $N$ the superlevel set
lies between the balls of squared radii $4[B_S-\delta]_+/k^2$ and
$4[B_S+\delta]_+/k^2$, the inner one empty when $B_S\le\delta$. The vectors $y$
form a translate of the lattice $(\sqrt L/N)\{z\in\Z^S:\sum_iz_i=0\}$ in $H_S$.
Its covolume is $\sqrt k(\sqrt L/N)^{k-1}$, since the basis $e_i-e_k$,
$1\le i<k$, has Gram determinant $k$. Its cells have diameter
$O(\sqrt L/N)\to0$, so a ball of bounded radius $\rho$ holds
$(|\mathbb B_{k-1}|\rho^{k-1}/\sqrt k+o(1))(N/\sqrt L)^{k-1}$ lattice points,
uniformly in the translate and in $\rho$. Letting $\delta\downarrow0$
gives~\eqref{eq:frontier-volume}, also when $B_S=0$. If $B_S<0$, the profile is
below 1 on the whole ball for large $N$, so there are no such bins.
\end{proof}

At zero field all supports of the same size have the same count, so the right
side of~\eqref{eq:frontier-volume} is multiplied by $\binom dk$ when we count all
bins with $k$ positive coordinates. For $k=2$, writing $y=(y_1,-y_1)$ gives the
profile $\exp(t-a_2-2y_1^2)$ and the count constant $\sqrt{2[t-a_2]_+}$, which
agrees with the binary window of~\cite{paperA} once the different centering is
taken into account. The covolume $\sqrt k$ is checked in Lean.

%======================================================================
% 5 (cont). The sentence in frontier_modes.tex that pointed to it
% Source: frontier_modes.tex, lines 90-93 of the version before 2026-09-26 trim.
%======================================================================
is $e^{\bar h_S+o(1)}$ near the center of the face. So the best bin on $S$
scores $\ell_S(t,h)$, and taking the $k$ largest fields gives the lines. The
proof is in Appendix~\ref{app:start-field}, which also counts the bins that
beat a vertex (Theorem~\ref{thm:frontier-superlevel}).

%======================================================================
% 6. Remark rem:tetra (the covariance)
% Source: app_start_field.tex, lines 347-364 of the version before 2026-09-26 trim.
%======================================================================
\subsection{The covariance}

\begin{remark}[the covariance]\label{rem:tetra}
Let the start be uniform, and let $K_{N,a}$ be the number of visits to state
$a$. The chain is stationary with transition matrix $\sigma I+(1-\sigma)\frac1dJ$,
whose $m$-th power is $\sigma^mI+(1-\sigma^m)\frac1dJ$. So the indicators of
state $a$ at time $s$ and of state $b$ at time $s+m$ have covariance
$\sigma^m(\frac{\delta_{ab}}d-\frac1{d^2})$. Summing over the 2 times and
evaluating the finite geometric sum gives, for $p<1$,
\[
 \Cov(K_{N,a},K_{N,b})=\Bigl(\frac{\delta_{ab}}d-\frac1{d^2}\Bigr)
 \Bigl(N\,\frac{1+\sigma}{1-\sigma}-\frac{2\sigma(1-\sigma^N)}{(1-\sigma)^2}\Bigr).
\]
Since $1-\sigma=dr$, for $1/d\le p<1$ the second term in the bracket is not
positive and $\Var(K_{N,a}/N)\le2/(d^2T)$. For $p<1/d$ the bracket is at most
$N+4$ while $T\le N$. So $\Var(K_{N,a}/N)=O(1/T)$ for every $p<1$, which is
what Proposition~\ref{prop:field-tilt}(3) uses.
\end{remark}

%======================================================================
% 7. Old moment bound in the proof of Lemma start-weights(2)
% Source: app_start_field.tex, lines 61-73 of the version before 2026-09-26 trim.
%======================================================================
It remains to bound the moments. Put $h_0=N/k$ and $x=h_0v$. The summands
of~\eqref{eq:simplex-positive-sum} are proportional to $w(m)$, and the proof of
Lemma~\ref{lem:simplex-integral} writes them as a constant times
$c(m)A(m)$. Here $c(m)=M!\prod_bx^{m_b}/(m_b!(m_b-1)!)$ are the weights of the
series $F_k(x)$, and $A(m)\in[0,1]$ with $1-A(m)\le M^2/h_0$. By
Lemma~\ref{lem:simplex-integral}, $\E_cM^2=\mathcal D^2F_k(x)/F_k(x)\le
C_2(1+x)^2$ for all $x>0$, with $C_2$ depending on $k$. Since $\sigma\ge\frac12$
we have $x=T/(k\sigma)\le2CL$, so $\E_cM^2\le h_0/2$ for $N\ge N_0$. Hence
\[
 \E_wM\le\E_wM^2\le\frac{\E_cM^2}{1-\E_cM^2/h_0}\le2C_2(1+x)^2.
\]
With $g\ge N/(2k)$ and $x\le2T$ this gives $\delta\le C_1(1+T)^2/N$.


%======================================================================
% 8. The remark after Lemma gaussian-phi-average on the case l=0
% Source: app_thresholds.tex, lines 102-104 of the version before 2026-09-26 trim.
%======================================================================
With $\ell=0$ the lemma is the Taylor step
$\frac1{\sqrt\pi}\int_{-w_0/2}^{w_0/2}e^{-y^2}(w_0+y)^q\,dy
=w_0^q\bigl(1+q(q-1)/(4w_0^2)+O(w_0^{-4})\bigr)$ of the Laplace expansions below.

%======================================================================
% 9. Lemma lem:simplex-integral and Lemma lem:simplex-laplace with their proofs
% Source: app_thresholds.tex, lines 135-220 of the version before 2026-09-26 trim.
%======================================================================
\begin{lemma}[positive integral and discrete comparison]\label{lem:simplex-integral}
For fixed $k\ge2$ define
\begin{equation}\label{eq:simplex-F}
F_k(x)=\sum_{m_1,\ldots,m_k\ge1}
M!\prod_{i=1}^k\frac{x^{m_i}}{m_i!(m_i-1)!}
=\int_0^\infty e^{-t}\bigl[\sqrt{xt}\,I_1(2\sqrt{xt})\bigr]^k\,dt,
\quad M=\sum_im_i.
\end{equation}
Let $n$ be a balanced positive count vector with total $N$, put $h=N/k$, and for
$0<u<1$ put $x=hu/(1-u)$. With $\mathcal D=x\,d/dx$,
\begin{equation}\label{eq:simplex-discrete-comparison}
0\le 1-\frac{S_n(u)}{(1-u)^{N-1}h^{1-k}F_k(x)/x}
\le\frac{\mathcal D^2F_k(x)}{hF_k(x)}
=O_k\Bigl(\frac{(1+x)^2}{h}\Bigr).
\end{equation}
In particular the relative error is $O_k((\log N)^2/N)$ when $x=O(\log N)$.
\end{lemma}

\begin{proof}
The equality in~\eqref{eq:simplex-F} follows from the series of $I_1$ and
Tonelli's theorem, using $M!=\int_0^\infty e^{-t}t^M\,dt$. The integral is finite
for every $x>0$, since $I_1(y)\le e^y$ reduces its tail to a polynomial times
$\exp(-t+2k\sqrt{xt})$. So the positive series and all its termwise derivatives
converge on compact subsets of $x>0$.

By Proposition~\ref{prop:dpmf}(c), $S_n$ is the positive
sum~\eqref{eq:simplex-positive-sum}. Define
\[
A_{i,m}=\frac{(m-1)!}{h^{m-1}}\binom{n_i-1}{m-1}
=\prod_{j=1}^{m-1}\Bigl(1-\frac{h-n_i+j}{h}\Bigr)_+ .
\]
Since $h-n_i>-1$, these are products of factors in $[0,1]$, also beyond their
finite support. Lemma~\ref{lem:clipped-product} gives
\[
0\le1-\prod_iA_{i,m_i}\le\frac1h\sum_i\frac{m_i(m_i+1)}2\le\frac{M^2}{h}.
\]
Replacing the binomial factors in~\eqref{eq:simplex-positive-sum} by
$h^{m_i-1}/(m_i-1)!$ gives the denominator
in~\eqref{eq:simplex-discrete-comparison}. Summing the deficit bound proves its
first inequality and the explicit upper bound.

For the quotient of derivatives with $x\ge1$ we use $M^2\le3x^2(1+1/x)^M$, which
holds since $\log(1+1/x)\ge1/(2x)$ and $\sup_{y\ge0}y^2e^{-y/2}=16/e^2<3$.
Positivity of the coefficients gives
\[
\frac{\mathcal D^2F_k(x)}{F_k(x)}\le3x^2\frac{F_k(x+1)}{F_k(x)}=O_k(x^2),
\]
where the last step follows from Lemma~\ref{lem:simplex-laplace} below. For
$0<x\le1$ the quotient extends continuously to $x=0$ with value $k^2$, since
$F_k(x)=k!x^k+O_k(x^{k+1})$.
\end{proof}

\begin{lemma}[Laplace expansion]\label{lem:simplex-laplace}
For each fixed $k\ge2$, as $x\to\infty$,
\begin{equation}\label{eq:simplex-laplace}
F_k(x)=\frac{k^{k/2+1}}{(4\pi)^{(k-1)/2}}\,x^{(k+1)/2}e^{k^2x}
\Bigl(1-\frac{k-1}{8kx}+O_k(x^{-2})\Bigr).
\end{equation}
\end{lemma}

\begin{proof}
In~\eqref{eq:simplex-F} put $t=w^2$ and $w_0=k\sqrt x$. The exponent becomes
$-w^2+2k\sqrt x\,w=k^2x-(w-w_0)^2$. On $|w-w_0|\le w_0/2$ we have
$2\sqrt x\,w\ge kx$, so the expansion of $I_1$ in
Section~\ref{sec:multistate-model} applies uniformly and gives
\[
\bigl[\sqrt x\,wI_1(2\sqrt x\,w)\bigr]^k
=\frac{x^{k/4}w^{k/2}e^{2k\sqrt x\,w}}{(4\pi)^{k/2}}
\Bigl(1-\frac{3k}{16\sqrt x\,w}+O_k(x^{-2})\Bigr).
\]
With $q=k/2+1$ and $w=w_0+y$, this range contributes
\[
\frac{2x^{k/4}e^{k^2x}}{(4\pi)^{k/2}}\int_{|y|\le w_0/2}e^{-y^2}
\Bigl(w^q-\frac{3k}{16\sqrt x}\,w^{q-1}+O_k(x^{-2})\,w^q\Bigr)dy .
\]
Lemma~\ref{lem:gaussian-phi-average} with $\ell=0$ evaluates the integrals of
$w^q$ and $w^{q-1}$. Since $q(q-1)/(4w_0^2)=(k+2)/(16kx)$ and
$3k/(16\sqrt x\,w_0)=3/(16x)$, the relative correction is
$(k+2)/(16kx)-3/(16x)=-(k-1)/(8kx)$. The leading term is
$2\sqrt\pi\,x^{k/4}w_0^q(4\pi)^{-k/2}e^{k^2x}
=k^{k/2+1}x^{(k+1)/2}(4\pi)^{-(k-1)/2}e^{k^2x}$. On the rest of $w>0$ the bound
$I_1(y)\le e^y$ makes the integrand at most
$2x^{k/2}w^{k+1}e^{k^2x-(w-w_0)^2}$, and $(w-w_0)^2\ge k^2x/4$ there. So this part
is at most $e^{7k^2x/8}$ times a power of $x$, which is exponentially smaller than
the main term.
\end{proof}

%======================================================================
% 10. Old first paragraph of the proof of the asymptotic assertions in Theorems simplex-window and simplex-crossing
% Source: app_thresholds.tex, lines 222-236 of the version before 2026-09-26 trim.
%======================================================================
\begin{proof}[Proof of the asymptotic assertions in
Theorems~\ref{thm:simplex-window} and~\ref{thm:simplex-crossing}]
Put $s=k-1$, let $n$ be balanced, and let $u=z/N$ with $c\log N\le z\le C\log N$.
Then $x=hu/(1-u)=z/(k(1-u))$, and Lemmas~\ref{lem:simplex-integral}
and~\ref{lem:simplex-laplace} give
\[
S_n(u)=(1-u)^{N-1}h^{-s}\frac{F_k(x)}x\Bigl(1+O\Bigl(\frac{z^2}N\Bigr)\Bigr),
\qquad
\frac{F_k(x)}x=\frac{k^{k/2+1}x^{s/2}e^{k^2x}}{(4\pi)^{s/2}}
\Bigl(1-\frac s{8kx}+O(x^{-2})\Bigr).
\]
Here $k^2x+(N-1)\log(1-u)-sz=O(z^2/N)$, $x^{s/2}=(z/k)^{s/2}(1+O(z/N))$ and
$s/(8kx)=s/(8z)+O(1/N)$. With $h^{-s}=k^sN^{-s}$ the powers of $k$ collect to
$k^{k+1/2}$, which gives~\eqref{eq:simplex-uniform-ratio}.


%======================================================================
% 10 (cont). Old first part of the proof of Theorem thm:simplex-interior
% Source: app_thresholds.tex, lines 263-305 of the version before 2026-09-26 trim.
%======================================================================
\begin{proof}[Proof of Theorem~\ref{thm:simplex-interior}]
The polynomial $S_n$ has nonnegative coefficients and lowest term $k!u^{k-1}$,
so $S_n(u)=1$ has a unique positive root. Put $v=u/(1-u)$, $y=Nv$ and
\[
G_\alpha(y)=\int_0^\infty e^{-t}\prod_i\bigl[\sqrt{\alpha_iyt}\,
I_1(2\sqrt{\alpha_iyt})\bigr]dt
=\sum_{m_i\ge1}M!\prod_i\frac{(\alpha_iy)^{m_i}}{m_i!(m_i-1)!}.
\]
Replacing each binomial coefficient in~\eqref{eq:simplex-positive-sum} by
$n_i^{m_i-1}/(m_i-1)!$ turns the sum into $G_\alpha(y)/(v\prod_in_i)$. Since
$(m-1)!\binom{n_i-1}{m-1}/n_i^{m-1}=\prod_{j<m}(1-j/n_i)_+$,
Lemma~\ref{lem:clipped-product} bounds the relative deficit of each term by
$\sum_im_i(m_i-1)/(2n_i)\le M^2/(2\varepsilon N)$. Let $\E_y$ be the mean under
weights proportional to the terms of $G_\alpha(y)$. As in the proof of
Lemma~\ref{lem:simplex-integral}, $\E_yM^2\le3y^2G_\alpha(y+1)/G_\alpha(y)$ for
$y\ge1$, and this is $O(y^2)$ uniformly by the expansion of $G_\alpha$ below. For
$y\le1$, $\E_yM^2$ is continuous in $(\alpha,y)$ and tends to $k^2$ as $y\to0$.
Hence
\[
S_n(u)=\frac{(1-u)^{N-1}}{v\prod_in_i}\,G_\alpha(y)
\Bigl(1+O_{k,\varepsilon}\Bigl(\frac{(1+y)^2}N\Bigr)\Bigr).
\]

Put $t=w^2$ and $w_0=A\sqrt y$, so the exponent is
$-w^2+2A\sqrt y\,w=A^2y-(w-w_0)^2$. On $|w-w_0|\le w_0/2$ every Bessel argument
has $2\sqrt{\alpha_iy}\,w\ge\sqrt\varepsilon\,y$, since $A\ge1$. So the expansion
of $I_1$ gives, uniformly,
\[
\prod_i\bigl[\sqrt{\alpha_iy}\,wI_1(2\sqrt{\alpha_iy}\,w)\bigr]
=\frac{(\prod_i\alpha_i)^{1/4}(\sqrt y\,w)^{k/2}e^{2A\sqrt y\,w}}{(4\pi)^{k/2}}
\bigl(1+O_{k,\varepsilon}(y^{-1})\bigr).
\]
Lemma~\ref{lem:gaussian-phi-average} with $\ell=0$ and $q=k/2+1$ evaluates the
remaining Gaussian integral, and the rest of $w>0$ is bounded as in the proof of
Lemma~\ref{lem:simplex-laplace}. So
\[
G_\alpha(y)=\frac{A^{k/2+1}(\prod_i\alpha_i)^{1/4}}{(4\pi)^{s/2}}\,
y^{(k+1)/2}e^{A^2y}\bigl(1+O_{k,\varepsilon}(y^{-1})\bigr).
\]
Substitute this and $v\prod_in_i=yN^s\prod_i\alpha_i$ into the previous display.
For $u=z/N$ in the stated range, $y=z+O(z^2/N)$ and
$(N-1)\log(1-u)=-z+O(z^2/N)$. This proves~\eqref{eq:simplex-interior-ratio}.


%======================================================================
% 10 (cont). Old sentence in simplex_thresholds.tex describing the asymptotic proof
% Source: simplex_thresholds.tex, lines 86-90 of the version before 2026-09-26 trim.
%======================================================================
which moves the root. For the asymptotic part, Proposition~\ref{prop:dpmf}(c)
writes $S_n$ as a positive sum. This sum is close to a Bessel integral $F_k$
(Lemma~\ref{lem:simplex-integral}), and the Laplace expansion of $F_k$
(Lemma~\ref{lem:simplex-laplace}) gives, for a balanced bin $n$ with $k$
positive coordinates and uniformly on $c\log N\le z\le C\log N$,

%======================================================================
% 11. Remark rem:last-exit-trees
% Source: app_thresholds.tex, lines 399-409 of the version before 2026-09-26 trim.
%======================================================================
\begin{remark}[last-exit trees]\label{rem:last-exit-trees}
In continuous time the trees appear directly. Huang, Kious, Sidoravicius and
Tarr\'es~\cite[Theorem~1]{huang2018} give the joint density of the local times,
the last-exit tree and the net crossing numbers of a walk with symmetric
conductances. Read on the slice of fixed total time and summed over the crossing
numbers by the generating function of the $I_m$~\cite[Eq.~10.35.1]{dlmf}, it is
an exact sum over spanning trees, each times a theta function. A Laplace step,
which we only sketch, turns each theta function into a common term proportional
to $\tree^{-1/2}$, and the trees then leave $\sqrt{\tree(\alpha)}$. Nothing in
this paper uses this remark.
\end{remark}

%======================================================================
% 12. Lemma lem:singleton-insertion and the sentence after it
% Source: app_boundary.tex, lines 353-372 of the version before 2026-09-26 trim.
%======================================================================
\begin{lemma}[exact insertion of a singleton state]
\label{lem:singleton-insertion}
Let $n$ have total size $M\ge1$ and let $\mathcal D=u\,d/du$. Then
\[
S_{(n,1)}(u)=\bigl[2u+(M-1)u^2\bigr]S_n(u)
+(u-u^2)\mathcal D S_n(u).
\]
\end{lemma}

\begin{proof}
Insert the new state into a word of length $M$ with $j$ changes. The 2 end
positions contribute $2u$, the $j$ change gaps contribute $ju$, and the
$M-1-j$ repeat gaps contribute $(M-1-j)u^2$. Removing the unique new letter
reverses this construction. Summing
$u^j[2u+(M-1)u^2+j(u-u^2)]$ over all words proves the formula.
\end{proof}

The singleton identity gives an exact finite-size check of the factor
$A^2z^2/N$ that one very small coordinate contributes in
\eqref{eq:simplex-boundary-profile}.

%======================================================================
% 13. Old motivation sentence and side-results remark in Subsection sec:vertex-completion
% Source: frontier_boundary.tex, lines 94-101 of the version before 2026-09-26 trim.
%======================================================================
Each rare run again needs about 2 switches and has about $b$ places, so the
rare counts enter through $bu^2$. For fixed rare counts,
$S_{(b,a)}(\sqrt{y/b})\to H(y)$ as $b\to\infty$
(Theorem~\ref{thm:simplex-fixed-vertex}). So the crossing is near
$\sqrt{y_a/b}$, and $Nu$ is of order $\sqrt N$ instead of $L$. The next
theorem shows that this leading term stays correct uniformly as long as
$|a|=o(b)$.


%======================================================================
% 13 (cont).
% Source: frontier_boundary.tex, lines 136-149 of the version before 2026-09-26 trim.
%======================================================================
\begin{remark}
Appendix~\ref{app:vertex} also proves 3 side results. For fixed rare counts,
Theorem~\ref{thm:simplex-fixed-vertex} gives the next term
$u_{b,a}=\sqrt{y_a/b}-(1+y_aQ(y_a)/H'(y_a))/b+O(b^{-3/2})$ with an explicit
polynomial $Q$, and for one rare count this is the 2-state fixed-edge
expansion of~\cite{paperA}. When $|a|\to\infty$ with $|a|=o(b)$ and
proportions $a_i/|a|$ bounded away from $0$,
Corollary~\ref{cor:vertex-shape-law} gives
$u_{b,a}\sim\ell\log|a|/(2A\sqrt{|a|b})$ with $A=\sum_i\sqrt{a_i/|a|}$.
Proposition~\ref{prop:vertex-probability-failure} shows that
$S_{(b,a)}(u)\sim H(bu^2)$ is not uniform over $|a|=o(b)$, while by
Theorem~\ref{thm:vertex-uniform-root} the root is.
\end{remark}


%======================================================================
% 14. Old opening of Appendix app:vertex
% Source: vertex_completion.tex, lines 1-6 of the version before 2026-09-26 trim.
%======================================================================
\section{Proofs near a vertex}
\label{app:vertex}

We use the notation of Subsection~\ref{sec:vertex-completion}. We first
prove Theorem~\ref{thm:vertex-uniform-root} and then the 3 side results
for fixed and for growing rare counts.

%======================================================================
% 14 (cont). Theorem thm:simplex-fixed-vertex, the all-singleton case, Corollary cor:vertex-shape-law and Proposition prop:vertex-probability-failure, with proofs
% Source: vertex_completion.tex, lines 156-326 of the version before 2026-09-26 trim.
%======================================================================
\begin{theorem}[a multistate Laguerre limit at a vertex]
\label{thm:simplex-fixed-vertex}
Fix $\ell\ge1$ and positive integers $a_1,\ldots,a_\ell$. Define
\[
Q(y)=\sum_{i<j}B'_{a_i}(y)B'_{a_j}(y)
\prod_{h\ne i,j}B_{a_h}(y).
\]
The sum defining $Q$ is zero when $\ell=1$. Uniformly on compact
subsets of $y>0$, as $b\to\infty$,
\begin{equation}\label{eq:simplex-vertex-ratio}
S_{(b,a)}\!\left(\sqrt{y/b}\right)
=H(y)+\frac{2\sqrt y}{\sqrt b}\bigl(H'(y)+yQ(y)\bigr)
+O_a(b^{-1}).
\end{equation}
Put $c_a=\sqrt{y_a}$. The crossing odds satisfy
\begin{equation}\label{eq:simplex-vertex-root}
u_{b,a}=\frac{c_a}{\sqrt b}
-\frac{1+y_aQ(y_a)/H'(y_a)}{b}
+O_a(b^{-3/2}).
\end{equation}
For one rare count the coefficient of $1/b$ is $-1$, as in the 2-state
fixed-edge expansion of~\cite{paperA}. For at least 2 rare counts it is
strictly less than $-1$.
\end{theorem}

\begin{proof}
We count words by their runs. Fix rare-run counts
$m_i\in\{1,\ldots,a_i\}$, put $m=\sum_im_i$, and let $q$ be the number of
runs of the large state. The rare run lengths can be chosen in
$\prod_i\binom{a_i-1}{m_i-1}$ ways. Successive large runs must be
separated, so $q\le m+1$, and the large run lengths have
$\binom{b-1}{q-1}$ choices. At $u=\sqrt{y/b}$, words with these run counts
contribute $O_a(b^{(q-m-1)/2})$. There are finitely many possible $m_i,q$,
independently of $b$, so only $q=m+1$ and $q=m$ matter through order
$b^{-1/2}$.

If $q=m+1$, the run word starts and ends in the large state and alternates
between the large state and rare states. There are $m!/\prod_im_i!$ rare
orders. Using $\binom{b-1}{m}=b^m/m!+O_a(b^{m-1})$ and summing gives
$H(y)$, with an $O_a(b^{-1})$ error.

If $q=m$, there are 2 possibilities. First, one endpoint is rare and the
run word otherwise alternates, which gives $2m!/\prod_im_i!$ orders. These
terms sum to $2\sqrt y\,H'(y)/\sqrt b$. Second, both endpoints are large
and exactly one pair of consecutive rare runs is present. The members of
that pair must be different states. For an ordered pair of states $i,j$,
contracting the marked adjacent pair shows that there are
$(m-1)!/[(m_i-1)!(m_j-1)!\prod_{h\ne i,j}m_h!]$ orders. Summing over the 2
orientations and all $i<j$ gives $2y^{3/2}Q(y)/\sqrt b$. Terms with
$q\le m-1$ and the next terms of the binomial expansions are $O_a(b^{-1})$,
uniformly for $y$ in a compact subset of $(0,\infty)$. This
proves~\eqref{eq:simplex-vertex-ratio}.

The polynomial $H$ has positive coefficients, so $H'(y_a)>0$. First
bracket the crossing using~\eqref{eq:simplex-vertex-ratio} to obtain
$bu_{b,a}^2\to y_a$. The finite run expansion may also be differentiated
uniformly on a compact neighborhood of $y_a$, and its derivative tends to
$H'(y_a)>0$. Then substitution of $u=c_a/\sqrt b+\theta/b$
into~\eqref{eq:simplex-vertex-ratio} forces
$\theta=-1-y_aQ(y_a)/H'(y_a)$, and the uniform remainder
gives~\eqref{eq:simplex-vertex-root}. If $\ell\ge2$, all terms of $Q(y)$
are positive for $y>0$, which proves the strict inequality.
\end{proof}

If all $\ell$ rare counts equal $1$, then $H(y)=y^\ell$ and $y_a=1$. This
gives the explicit case
\[
u_{b,1,\ldots,1}=b^{-1/2}-\frac{\ell+1}{2b}
+O_\ell(b^{-3/2}).
\]
An independent exact formula in this case is
\[
S_{(b,1,\ldots,1)}(u)=\ell!
\sum_{q=1}^{\min(b,\ell+1)}
\binom{\ell+1}{q}\binom{b-1}{q-1}u^{q+\ell-1}.
\]
To see this, order the $\ell$ distinct rare letters, choose the $q$ gaps
among their $\ell+1$ end and interior gaps that contain large runs, and
choose the positive lengths of those runs. In particular 2 singleton rare
states give $u_{b,1,1}=b^{-1/2}-\tfrac32b^{-1}+O(b^{-3/2})$. This agrees
with the exact polynomial
$S_{(b,1,1)}(u)=6u^2+6(b-1)u^3+(b-1)(b-2)u^4$.
Theorem~\ref{thm:simplex-fixed-vertex} is not uniform in the rare counts.
For growing rare counts, Theorem~\ref{thm:vertex-uniform-root} gives the
leading term uniformly for $|a|=o(b)$, and the next result gives its shape.

\begin{corollary}[an explicit shape law for growing rare counts]
\label{cor:vertex-shape-law}
Suppose $|a|\to\infty$, $|a|=o(b)$, and the actual proportions
$\alpha_i=a_i/|a|$ stay in a fixed compact subset of the positive
$\ell$-simplex. Put $A=\sum_i\sqrt{\alpha_i}$. Then
\[
 u_{b,a}\sim\frac{\ell\log|a|}{2A\sqrt{|a|b}}.
\]
In more detail, the root $y_a$ satisfies
\begin{align*}
 2A\sqrt{|a|y_a}
 ={}&\ell\log|a|-\frac\ell2\log\log|a|
 +\frac\ell2\log\!\left(\frac{8\pi A}{\ell}\right)
 +\frac34\sum_i\log\alpha_i
 +O_\ell\!\left(\frac{\log\log|a|}{\log|a|}
                  +\frac{(\log|a|)^2}{|a|}\right).
\end{align*}
Replacing $\sqrt{|a|y_a}$ on the left by $\sqrt{|a|b}\,u_{b,a}$ adds an
error $O_\ell(\log|a|\sqrt{|a|/b})$.
\end{corollary}

\begin{proof}
Lemma~\ref{lem:uniform-laguerre-bessel} gives
$H(y)=y^\ell\prod_i\Phi(a_iy)(1+O_\ell(y))$, uniformly as $y\to0$.
Write $x=\sqrt{|a|y}$. Bracketing first at fixed small and large
multiples of $\log|a|$ shows that the root has $x\asymp\log|a|$.
The expansion of $I_1$ in Section~\ref{sec:multistate-model} then gives,
uniformly over the stated proportions,
\[
 H(x^2/|a|)=
 \frac{x^{\ell/2}e^{2Ax}}
 {|a|^\ell(4\pi)^{\ell/2}(\prod_i\alpha_i)^{3/4}}
 \left(1+O_\ell(x^{-1}+x^2/|a|)\right).
\]
Take logarithms, substitute $x=(\ell/(2A))\log|a|+O_\ell(\log\log|a|)$ in
the term $(\ell/2)\log x$, and use Theorem~\ref{thm:vertex-uniform-root}.
This proves the claims. The constants also depend on the fixed compact set
of proportions.
\end{proof}

\begin{proposition}[probability approximation can fail while roots agree]
\label{prop:vertex-probability-failure}
Let $a=\lfloor b/\log b\rfloor$, let $y_a$ solve $B_a(y_a)=1$,
and set $u_0=\sqrt{y_a/b}$. Then
\[
 H(y_a)=B_a(y_a)^2=1,\qquad
 S_{(b,a,a)}(u_0)\longrightarrow\infty,
 \qquad \frac{u_{b,a,a}}{u_0}\longrightarrow1.
\]
Hence the crossing can be approximated uniformly even where
$S(u)\sim H(bu^2)$ fails, so that approximation is not uniform over
$|a|=o(b)$.
\end{proposition}

\begin{proof}
We keep only the run words that start and end in the large state and have
exactly one adjacent pair of distinct rare states. The run count in the
proof of Theorem~\ref{thm:simplex-fixed-vertex} shows that their
contribution is the polynomial $2uyB_a'(y)^2$, with each term of total
rare-run count $m$ also multiplied by $\prod_{j=1}^{m-1}(1-j/b)_+$.
Give $J$ probabilities proportional to $\binom{a-1}{j}y^j/j!$, the
coefficients of $B_a'(y)$. Its differential equation
$y(B_a')''+(1+y)(B_a')'-(a-1)B_a'=0$ implies
$\E J^2+y\E J=(a-1)y$, and hence
$\E J\le\sqrt{ay}$ and $\E J^2\le ay$.
For 2 independent copies, $m=J_1+J_2+2$, so
\[
 \E[m(m-1)]\le4ay+6\sqrt{ay}+2.
\]
Lemma~\ref{lem:clipped-product} gives
$1-\prod_{j=1}^{m-1}(1-j/b)_+\le m(m-1)/(2b)$, and therefore
\[
 S_{(b,a,a)}(u_0)\ge
 \frac{2\kappa_a(y_a)^2}{bu_0}
 \left(1-\frac{4ay_a+6\sqrt{ay_a}+2}{2b}\right),
\]
where $\kappa_a=yB_a'/B_a$ as in the proof of
Lemma~\ref{lem:vertex-elasticity}. Corollary~\ref{cor:vertex-shape-law}
with $\ell=1$ gives $\sqrt{ay_a}\sim\tfrac12\log a$, and the bounds on
$\kappa_a$ in the proof of Lemma~\ref{lem:vertex-elasticity} give
$\kappa_a(y_a)\sim\sqrt{ay_a}$. So the lower bound is asymptotic to
$\sqrt{a/b}\log a\sim\sqrt{\log b}$, which diverges, although slowly.
The root assertion follows from Theorem~\ref{thm:vertex-uniform-root},
since $2a=o(b)$.
\end{proof}

%======================================================================
% 15. Main-text sentences that pointed to removed results (introduction, related work, open questions)
% Source: problem131_multistate.tex, lines 141-146 of the version before 2026-09-26 trim.
%======================================================================
(Theorem~\ref{thm:simplex-boundary-crossover}). Near a vertex, where one count
$b$ is large and the other counts $a_1,\ldots,a_\ell$ have total $|a|=o(b)$, the
crossing has $u=r/p=\sqrt{y_a/b}\,(1+O(\sqrt{(|a|+1)/b}))$, where $y_a$ is the
root of a
product of Laguerre-type polynomials (Theorem~\ref{thm:vertex-uniform-root}). This holds
uniformly even where the matching approximation to the probabilities fails.

%======================================================================
% 15 (cont).
% Source: problem131_multistate.tex, lines 209-211 of the version before 2026-09-26 trim.
%======================================================================
Huang, Kious, Sidoravicius and Tarr\'es~\cite{huang2018} give a Bessel formula
for the density of local times on finite graphs, and
Remark~\ref{rem:last-exit-trees} relates it to our tree sums.

%======================================================================
% 15 (cont).
% Source: problem131_multistate.tex, lines 279-281 of the version before 2026-09-26 trim.
%======================================================================
Paper~C. Example~\ref{ex:simplex-edge-mode} and
Proposition~\ref{prop:vertex-probability-failure} show that the simplest
strengthenings of our results are false.
